\section{Implications for individuals, organisations, and society}
\label{sec:implications}

The measurements in this paper are about balance sheets, but the exposure they describe is
ultimately borne by people, and the choices it forces are management and policy choices rather
than accounting ones. This section sets out what follows for four audiences.

\paragraph{For the people displacement reaches.} The household result that should travel furthest
is that the thinnest buffers are not at the bottom of the wage distribution. They are one quintile
up, among households that have taken on the commitments of the households above them without the
income or the assets to carry them: a median of \result{QTwoRunway} months of liquid runway, with
\result{QTwoUnderOne} percent holding less than a single month. A relief programme designed around
the bottom quintile misses them. So does a programme designed around occupation, because the
difference between physically and cognitively exposed households is almost entirely a difference
in pay, and what predicts whether a household survives a lost paycheque is what it earns rather
than what kind of work produced the earnings.

Two further findings bear directly on who carries the cost. If automation arrives through hiring
that does not happen rather than through layoffs, the private loss shifts onto people at the start
of their working lives, who hold student debt and little else, and it becomes invisible to the
separations data that most early-warning systems watch. And a claim that is serviced from wages is
a claim on somebody's pay: the reason this paper insists on the phrase \emph{labor backing} rather
than some more neutral term is that the aggregate is made of individual obligations that do not
adjust when the income behind them stops.

\paragraph{For managers and organisations.} The useful reframing for an institution is that wage
risk is not a credit category to be underwritten but a concentration to be governed. Three things
follow from the measured results. First, the exposure is uneven in a way the aggregate hides:
lenders whose books are unsecured consumer credit are the first to break and mortgage portfolio
lenders are among the last, which is close to the reverse of the intuition that housing is where
household risk lives. Second, an institution cannot protect itself against this exposure by
tightening household underwriting, because most of what reaches it arrives through the second
round, the spending and house prices and business credit that follow other people's lost wages,
and not through its own borrowers defaulting. Third, an institution that holds both sides of the
bet, insurers and private credit funds in particular, should report them together rather than in
separate parts of a disclosure, because the correlation between them is the whole point.

For a management reader the governance question is therefore not whether to buy a model of
automation risk. It is whether the institution can state, in a form a board can read, how much of
its income and how much of its collateral depends on wages continuing to be paid, and what it has
agreed in advance to do if they are not. Every contract instrument in this paper works only if it
was written before the event.

\paragraph{Strategic implications.} The strategic finding is about ownership, and it cuts against
the reflex to reach for a tax. The state holds the wage side of a two-sided bet and almost none of
the other side, and the only claim it has on the other side is a tax that reaches
\result{TauKBarkai} percent of the surplus against the \result{RequiredEasy} to
\result{RequiredHard} percent that replacing the lost wage taxes would need. Closing that position
by buying into the asset is priced in Section~\ref{sec:institutions} and it is expensive:
\result{StakeForTenPctBn} billion dollars to close about \result{GapClosedAtTenPct} percent of the
gap. Neither instrument is comfortable, and a government that notices this late will find that the
price of the asset has risen for the same reason its revenue fell.

The strategic reading for an organisation is parallel. An exposure that everyone holds cannot be
diversified away inside the system: our measurement is that it is spread almost evenly, so the
system in aggregate cannot rotate out of it even though any single institution can. That makes it
a candidate for pre-agreed rules rather than for portfolio management, and it makes disclosure the
instrument with the best ratio of value to cost.

\paragraph{For society and fiscal equity.} Three points are worth stating plainly. The first is
that the claims wages pay are spread far more evenly than the wages themselves, so a proportional
fall in pay is regressive before any policy response is chosen. The second is that the same
ownership concentration runs through both halves of this paper: the reason the capital tax reaches
so little is that most corporate claims are held in forms the individual income tax does not
reach, and the reason an AI bust transmits so weakly to household spending is that the assets that
would fall are held by the households least likely to change their spending. Concentrated
ownership makes the upside hard to tax and the downside narrow in its incidence, and those are the
same fact seen from two directions.

The third is about what a measurement like this is for. We deliberately do not use it to argue for
screening individual borrowers by occupation: Section~\ref{sec:institutions} sets out why the
evidence does not support it and why the legal exposure runs through discriminatory effects. A
statistic that identifies where wage risk sits should be used to size public instruments and to
make exposures visible, not to price individuals out of credit on the basis of what they do for a
living. That is a choice about how the measurement is used, and it belongs in the paper that
produces it.
